% Source PDF pages 48--52; manual pages 2-19--2-23.
\subsection{Body-Fixed Cartesian Coordinate System}\label{sec:body-fixed}

A body-fixed coordinate system is commonly used to define components of the
vehicle's aerodynamic and propulsive forces. For example, the aerodynamic
coefficients $C_A$ and $C_N$ define one force component along the vehicle
centerline and another normal to it.

The origin of the system is at the vehicle's center of mass. Its $x$ axis is the
vehicle longitudinal axis, the $y$ axis points toward the right wing, and the $z$
axis points toward the bottom of the vehicle, as shown in
\cref{fig:body-fixed-coordinates}. The unit vectors
$\hat{x}_b$, $\hat{y}_b$, and $\hat{z}_b$ are aligned with these axes.

\begin{figure}[htbp]
  \centering
  \resizebox{0.58\linewidth}{!}{\begin{tikzpicture}[scale=0.82,>=Latex,line cap=round,line join=round]
  \coordinate (cg) at (0,0);
  % Simplified vehicle planform.
  \draw[line width=0.9pt] (-3.9,1.25) .. controls (-3.4,1.55) and (-2.0,1.10) .. (-0.3,0.20)
    -- (2.8,-1.10) .. controls (3.20,-1.30) and (3.45,-0.85) .. (3.10,-0.55)
    -- (0.45,0.35) .. controls (-1.60,1.55) and (-3.50,1.75) .. cycle;
  \draw[line width=0.9pt] (-0.55,0.20) -- (-0.20,-1.95) -- (1.40,-1.15) -- (0.20,0.20);
  \draw[line width=0.9pt] (0.25,0.35) -- (2.25,1.15) -- (2.95,0.25) -- (1.15,-0.15);
  \draw[line width=0.9pt] (1.35,0.05) -- (1.85,1.55) -- (2.15,1.45) -- (2.45,-0.35);
  \draw[line width=0.9pt] (3.05,-0.55) ellipse (0.32 and 0.46);
  \fill (cg) circle (2pt) node[above left] {cg};
  \draw[->,line width=1.1pt] (cg) -- (-2.25,1.02) node[above left] {$\hat{x}_b$};
  \draw[->,line width=1.1pt] (cg) -- (1.42,0.75) node[above right] {$\hat{y}_b$};
  \draw[->,line width=1.1pt] (cg) -- (0,-1.45) node[below] {$\hat{z}_b$};
\end{tikzpicture}
}
  \caption{Body-fixed coordinates.}
  \label{fig:body-fixed-coordinates}
\end{figure}

The body system is related to the local geodetic horizon system by three geodetic
Euler angles: yaw $\Psi_{gd}$, pitch $\Theta_{gd}$, and roll $\Phi_{gd}$. The yaw
and pitch angles are analogous to the horizontal and vertical flight-path angles;
the difference is that the vehicle longitudinal axis replaces the velocity vector.
The rotation sequence from the local geodetic horizon system to the body system is
shown in \cref{fig:geodetic-euler-angles}.

\begin{figure}[htbp]
  \centering
  \resizebox{0.90\linewidth}{!}{\begin{tikzpicture}[>=Latex,line cap=round,line join=round,font=\sffamily\scriptsize]
  % Panel 1: yaw about z_gd.
  \begin{scope}[shift={(0,0)}]
    \coordinate (O) at (0,0);
    \draw[->,line width=0.9pt] (O)--(0,1.55) node[above] {$\hat{x}_{gd}$};
    \draw[->,line width=0.9pt] (O)--(1.55,0) node[right] {$\hat{y}_{gd}$};
    \node[left] at (-0.18,0.05) {$\hat{z}_{gd}$};
    \node at (-0.15,-0.26) {$\otimes$};
    \draw[->,line width=1.05pt] (O)--(0.88,1.12) node[above right] {$\hat{x}_{1}$};
    \draw[->,line width=1.05pt] (O)--(1.12,-0.88) node[below right] {$\hat{y}_{1}$};
    \draw[->] (0,0.68) arc[start angle=90,end angle=51,radius=0.68];
    \node at (0.42,0.78) {$\Psi_{gd}$};
    \node[align=center] at (0.72,-1.55) {Yaw rotation\\about $\hat{z}_{gd}$};
  \end{scope}

  % Panel 2: pitch about y_1; the y axis is normal to the page.
  \begin{scope}[shift={(4.0,0)}]
    \coordinate (O) at (0,0);
    \node[align=center] at (0,1.72) {$\hat{y}_{1}$ is into\\the page};
    \node at (0,1.18) {$\otimes$};
    \draw[->,line width=0.9pt] (O)--(-0.85,1.15) node[above left] {$\hat{x}_{1}$};
    \draw[->,line width=0.9pt] (O)--(-0.78,-1.18) node[left] {$\hat{z}_{1}$};
    \draw[->,line width=1.05pt] (O)--(1.50,0.20) node[right] {$\hat{x}_{2}$};
    \draw[->,line width=1.05pt] (O)--(-0.20,-1.48) node[right] {$\hat{z}_{2}$};
    \draw[->] (-0.50,0.66) arc[start angle=127,end angle=8,radius=0.80];
    \node at (0.45,0.72) {$\Theta_{gd}$};
    \node[align=center] at (0.25,-2.05) {Pitch rotation\\about $\hat{y}_{1}$};
  \end{scope}

  % Panel 3: roll about x_2.
  \begin{scope}[shift={(8.1,0)}]
    \coordinate (O) at (0,0);
    \node[align=center] at (0,1.72) {$\hat{x}_{2}$ is into\\the page};
    \node at (0,1.18) {$\otimes$};
    \draw[->,line width=0.9pt] (O)--(1.45,0) node[right] {$\hat{y}_{2}$};
    \draw[->,line width=0.9pt] (O)--(0,-1.45) node[left] {$\hat{z}_{2}$};
    \draw[->,line width=1.05pt] (O)--(1.12,0.82) node[above right] {$\hat{y}_{b}$};
    \draw[->,line width=1.05pt] (O)--(0.82,-1.12) node[right] {$\hat{z}_{b}$};
    \draw[->] (0.74,0) arc[start angle=0,end angle=36,radius=0.74];
    \node at (0.92,0.30) {$\Phi_{gd}$};
    \node[align=center] at (0.55,-2.05) {Roll rotation\\about $\hat{x}_{2}$};
  \end{scope}
\end{tikzpicture}
}
  \caption{Geodetic yaw, pitch, and roll angles.}
  \label{fig:geodetic-euler-angles}
\end{figure}

The first rotation is a yaw about $\hat{z}_{gd}$, producing a temporary system
$\hat{x}_1$, $\hat{y}_1$, and $\hat{z}_1$, with
$\hat{z}_1=\hat{z}_{gd}$. The second rotation is a pitch about $\hat{y}_1$,
producing $\hat{x}_2$, $\hat{y}_2$, and $\hat{z}_2$, with
$\hat{y}_2=\hat{y}_1$. The final rotation is a roll about $\hat{x}_2$, and
$\hat{x}_2=\hat{x}_b$.

The same sequence relates the body system to the local geocentric horizon system
through the geocentric Euler angles $\Psi_{gc}$, $\Theta_{gc}$, and
$\Phi_{gc}$, and to the inertial-platform system through the platform Euler angles
$\Psi_p$, $\Theta_p$, and $\Phi_p$. In those cases the initial system in
\cref{fig:geodetic-euler-angles} is replaced by the local geocentric or inertial
platform system. For the geodetic case, the rotations are
\begin{equation}
  \begin{bmatrix}
    \hat{x}_b\\[0.2em]
    \hat{y}_b\\[0.2em]
    \hat{z}_b
  \end{bmatrix}
  =
  \TaosRollMatrix{\Phi_{gd}}
  \TaosPitchMatrix{\Theta_{gd}}
  \TaosYawMatrix{\Psi_{gd}}
  \begin{bmatrix}
    \hat{x}_{gd}\\[0.2em]
    \hat{y}_{gd}\\[0.2em]
    \hat{z}_{gd}
  \end{bmatrix}.
  \label{eq:body-euler-transformation}
\end{equation}
When expanded, \cref{eq:body-euler-transformation} gives
\begin{equation}
  \hat{x}_b
  = \cos\Theta_{gd}\cos\Psi_{gd}\,\hat{x}_{gd}
  + \cos\Theta_{gd}\sin\Psi_{gd}\,\hat{y}_{gd}
  - \sin\Theta_{gd}\,\hat{z}_{gd},
  \label{eq:body-x-from-geodetic-euler}
\end{equation}
\begin{equation}
  \begin{aligned}
  \hat{y}_b
  ={}&
  (\sin\Phi_{gd}\sin\Theta_{gd}\cos\Psi_{gd}
   -\cos\Phi_{gd}\sin\Psi_{gd})\,\hat{x}_{gd}\\
  &+(\sin\Phi_{gd}\sin\Theta_{gd}\sin\Psi_{gd}
   +\cos\Phi_{gd}\cos\Psi_{gd})\,\hat{y}_{gd}\\
  &+\sin\Phi_{gd}\cos\Theta_{gd}\,\hat{z}_{gd},
  \end{aligned}
  \label{eq:body-y-from-geodetic-euler}
\end{equation}
\begin{equation}
  \begin{aligned}
  \hat{z}_b
  ={}&
  (\cos\Phi_{gd}\sin\Theta_{gd}\cos\Psi_{gd}
   +\sin\Phi_{gd}\sin\Psi_{gd})\,\hat{x}_{gd}\\
  &+(\cos\Phi_{gd}\sin\Theta_{gd}\sin\Psi_{gd}
   -\sin\Phi_{gd}\cos\Psi_{gd})\,\hat{y}_{gd}\\
  &+\cos\Phi_{gd}\cos\Theta_{gd}\,\hat{z}_{gd}.
  \end{aligned}
  \label{eq:body-z-from-geodetic-euler}
\end{equation}
The function \texttt{euler\_to\_body} computes the body unit vectors from
geodetic, geocentric, or inertial-platform Euler angles. These equations assume
that the Euler angles are known. In many guidance problems the body attitude is
instead specified through aerodynamic angles, which relate the body and wind
systems as described in \cref{sec:wind-coordinate}.

When the body unit vectors are known from aerodynamic angles, the corresponding
Euler angles are recovered from the unit vectors. For the geodetic angles,
\begin{equation}
  \cos\Psi_{gd}
  = \frac{\hat{x}_b\mathbin{\cdot}\hat{x}_{gd}}
  {\sqrt{(\hat{x}_b\mathbin{\cdot}\hat{x}_{gd})^2
       +(\hat{x}_b\mathbin{\cdot}\hat{y}_{gd})^2}},
  \label{eq:geodetic-yaw-cosine}
\end{equation}
\begin{equation}
  \sin\Psi_{gd}
  = \frac{\hat{x}_b\mathbin{\cdot}\hat{y}_{gd}}
  {\sqrt{(\hat{x}_b\mathbin{\cdot}\hat{x}_{gd})^2
       +(\hat{x}_b\mathbin{\cdot}\hat{y}_{gd})^2}},
  \label{eq:geodetic-yaw-sine}
\end{equation}
\begin{equation}
  \cos\left(\Theta_{gd}+\frac{\pi}{2}\right)
  = \hat{x}_b\mathbin{\cdot}\hat{z}_{gd},
  \label{eq:geodetic-pitch-from-body-axis}
\end{equation}
\begin{equation}
  \cos\Phi_{gd}=\hat{y}_1\mathbin{\cdot}\hat{y}_b.
  \label{eq:geodetic-roll-from-body-axis}
\end{equation}
The yaw angle is obtained from the projection of $\hat{x}_b$ onto the local
geodetic horizon plane. The pitch angle is obtained from the projection of
$\hat{x}_b$ onto $\hat{z}_{gd}$. The roll angle is obtained by projecting
$\hat{y}_b$ onto the temporary $\hat{y}_1$ axis. This axis is perpendicular to
$\hat{x}_1$ and $\hat{z}_{gd}$:
\begin{equation}
  \hat{y}_1=\hat{z}_{gd}\times\hat{x}_1.
  \label{eq:temporary-yaw-y-axis}
\end{equation}
The temporary $\hat{x}_1$ axis satisfies
\begin{equation}
  \begin{bmatrix}
    \hat{x}_{gd}\\[0.2em]
    \hat{y}_{gd}\\[0.2em]
    \hat{z}_{gd}
  \end{bmatrix}
  \mathbin{\cdot}\hat{x}_1
  =
  \begin{bmatrix}
    \cos\Psi_{gd}\\[0.2em]
    \sin\Psi_{gd}\\[0.2em]
    0
  \end{bmatrix},
  \label{eq:temporary-yaw-x-linear-system}
\end{equation}
where the unit-vector components form the rows of an orthogonal matrix. Its inverse
is its transpose, so
\begin{equation}
  \hat{x}_1
  =
  \begin{bmatrix}
    \hat{x}_{gd}&\hat{y}_{gd}&\hat{z}_{gd}
  \end{bmatrix}
  \begin{bmatrix}
    \cos\Psi_{gd}\\[0.2em]
    \sin\Psi_{gd}\\[0.2em]
    0
  \end{bmatrix}.
  \label{eq:temporary-yaw-x-axis}
\end{equation}
This method is singular when the body points straight up or down, so that
$\hat{x}_b$ is aligned with $\hat{z}_{gd}$ and yaw is undefined. TAOS then sets
\begin{equation}
  \Psi_{gd}=\Phi_{gd}=0,
  \qquad
  \Theta_{gd}=\frac{\pi}{2}
    (\hat{x}_b\mathbin{\cdot}\hat{z}_{gd}).
  \label{eq:euler-angle-pole-convention}
\end{equation}
The geocentric and inertial-platform Euler angles are calculated by the same
procedure after substituting the corresponding reference-frame unit vectors. These
calculations are performed by \texttt{euler\_angles}.

\taossubhead{Euler-angle output variables}

\begin{longtable}{@{}>{\ttfamily}p{0.17\textwidth}p{0.18\textwidth}p{0.57\textwidth}@{}}
  yawgc   & $\Psi_{gc}$ & Geocentric Euler yaw angle.\\
  yawgd   & $\Psi_{gd}$ & Geodetic Euler yaw angle.\\
  yawi    & $\Psi_i$    & Inertial-platform Euler yaw angle.\\
  pitchgc & $\Theta_{gc}$ & Geocentric Euler pitch angle.\\
  pitchgd & $\Theta_{gd}$ & Geodetic Euler pitch angle.\\
  pitchi  & $\Theta_i$    & Inertial-platform Euler pitch angle.\\
  rollgc  & $\Phi_{gc}$ & Geocentric Euler roll angle.\\
  rollgd  & $\Phi_{gd}$ & Geodetic Euler roll angle.\\
  rolli   & $\Phi_i$    & Inertial-platform Euler roll angle.\\
\end{longtable}

\taossubhead{Body-coordinate vector components}

\begin{longtable}{@{}>{\ttfamily}p{0.17\textwidth}p{0.22\textwidth}p{0.53\textwidth}@{}}
  accebx & $\vec a_{\oplus}\mathbin{\cdot}\hat{x}_b$ &
    Earth-fixed acceleration component in the body $x$ direction.\\
  acceby & $\vec a_{\oplus}\mathbin{\cdot}\hat{y}_b$ &
    Earth-fixed acceleration component in the body $y$ direction.\\
  accebz & $\vec a_{\oplus}\mathbin{\cdot}\hat{z}_b$ &
    Earth-fixed acceleration component in the body $z$ direction.\\
  accibx & $\vec a_I\mathbin{\cdot}\hat{x}_b$ &
    Inertial acceleration component in the body $x$ direction.\\
  acciby & $\vec a_I\mathbin{\cdot}\hat{y}_b$ &
    Inertial acceleration component in the body $y$ direction.\\
  accibz & $\vec a_I\mathbin{\cdot}\hat{z}_b$ &
    Inertial acceleration component in the body $z$ direction.\\
  velebx & $\vec V_{\oplus}\mathbin{\cdot}\hat{x}_b$ &
    Earth-relative velocity component in the body $x$ direction.\\
  veleby & $\vec V_{\oplus}\mathbin{\cdot}\hat{y}_b$ &
    Earth-relative velocity component in the body $y$ direction.\\
  velebz & $\vec V_{\oplus}\mathbin{\cdot}\hat{z}_b$ &
    Earth-relative velocity component in the body $z$ direction.\\
  velibx & $\vec V_I\mathbin{\cdot}\hat{x}_b$ &
    Inertial velocity component in the body $x$ direction.\\
  veliby & $\vec V_I\mathbin{\cdot}\hat{y}_b$ &
    Inertial velocity component in the body $y$ direction.\\
  velibz & $\vec V_I\mathbin{\cdot}\hat{z}_b$ &
    Inertial velocity component in the body $z$ direction.\\
\end{longtable}
